\section{Some Elementary Results}
In this section we give some elementary results on the number of solutions. We present some classical theorems such as those of K\"{o}nig-Rados, Chevalley, and Warning. We also establish elementary upper bounds for the number of solutions and results on the expected order of magnitude.
\subsection{Univariate Polynomials}
We start of with univariate polynomials. Let $f\in\bbF_q[X]$ is an univariate polynomial Naturally the number of solutions in $\bbF_q$ for the equation $f(X)=0$ is basically the degree of $\gcd(f(X),X^q-X)$. Since checking if $0$ is a solution or not is very easy to check we'll only consider non-zero solutions. So the number of non-zero solutions of $f(X)=0$ is basically the degree of $\gcd(f(X),X^{q-1}-1)$. Therefore we may always assume without loss of generality that $\deg(f)\leq q-1$. Suppose $$f(X)=a_{q-1}\cdot X^{q-1}+a_{q-2}\cdot X^{q-2}+\cdots+a_1\cdot X+a_0$$where $a_i\in\bbF_q$ for all $0\leq i\leq q-1$. Since for all $\alpha\in\bbF_q^*$ we have $\alpha^{q-1}=1$ the variable $X^{q-1}$ will always evaluate to be $1$ for all non-zero values of $X$. Therefore the number of solutions of $f$ is same as of the equation $a_{q-2}\cdot X^{q-2}+\cdots+a_1\cdot X+(a_0+a_{q-1})=0$. Therefore we can always assume $\deg(f)\leq q-2$. So suppose we have a polynomial $f\in\bbF_q[X]$ such that $$f(X)=a_{q-2}\cdot X^{q-2}+\cdots+a_1\cdot X+a_0$$where $a_i\in\bbF_q$ for all $0\leq i\leq q-2$. 

Now we can also analyze the number of solutions of $f$ using matrix methods. From $f$ we can define the following matrix $$M_f=\mat{a_0 & a_1 & \cdots & a_{q-2}\\ a_1 & a_2 & \cdots & a_0\\ \vdots & \vdots & \ddots & \vdots\\ a_{q-2} & a_0 & \cdots & a_{q-3}}$$This matrix is a \emph{left circulant matrix} i.e. each row is obtained from the preceding row by a cyclic shift of the entries to the left. The number of non-zero solutions of $f(X)=0$ and the rank of the matrix $M_f$ are very related. 
\begin{Theorem}{K\"{o}nig-Rados Theorem}{thm:konig-rados}
  Let $f(X)=\sum_{i=0}^{q-2}a_i\cdot X^i\in\bbF_q[X]$ be an univariate polynomial. Then $$\ov{Z}(f)=q-1-\Rank(M_f)$$  
\end{Theorem}
\begin{proof}
  Since $\deg(f)\leq q-2$ we have $\ov{Z}(f)\leq q-2$. $\ov{Z}(f)=q-1$ rises only when $f$ is a zero polynomial i.e. $M_f$ is a zero matrix and in that case $\Rank(M_f)=0$ which satisfies the theorem. Let $\gm$ be the primitive element of $\bbF_q$. Then consider the \emph{Vandermonde Matrix}, $V(\gm,\gm^2,\dots, \gm^{q-1})$: $$V=\mat{1&1&\cdots&1\\ \gm & \gm^2 &\cdots & \gm^{q-1}\\ \vdots & \vdots & \ddots & \vdots\\ \gm^{q-2} & \gm^{2 (q-2)} & \cdots &\gm^{(q-1)(q-2)}}$$Then using $\gm^{i\cdot (q-1)}=1$ for all $i\in[q-1]$ we obtain $$M_f\cdot V=\mat{f(\gm) & f(\gm^2) & \cdots & f(\gm^{q-1})\\ \gm^{-1}f(\gm) & \gm^{-2} f(\gm^2) & \cdots & \gm^{-(q-1)}f(\gm^{q-1})\\ \vdots & \vdots & \ddots & \vdots \\ \gm^{-(q-2)}f(\gm) & \gm^{-2(q-2)}f(\gm^2) & \cdots & \gm^{-(q-1)(q-2)}f(\gm^{q-1})}$$Now for any $i\in[q-1]$ we have $f(\gm^i)=0$ if and only if  the $i^{th}$ row becomes zero. Therefore the $\Rank(M_f\cdot V)=q-1-\ov{Z}(f)$. Since $V$ is the Vandermonde matrix $V$ and $M_f\cdot V$ have the same rank. So the theorem follows. 
\end{proof}

In the next subsections we will give simple upper and lower bounds for solutions multivariate polynomial equations or system of polynomial equations. 
\subsection{Chevalley--Warning Theorem}
We begin with two elementary lemmas about power sums that will serve as the key tools.
\begin{lemma}{}{lem:power-sum}
  Let $k$ be a non-negative integer. Suppose for $k=0$, $0^0=1$. Then $$\sum_{\alpha\in\bbF_q}\alpha^k=\begin{cases}
    0 & \text{if $k=0$ or $q-1\nmid k$}\\ 
    -1 & \text{if $q-1\mid k$}
  \end{cases}$$
\end{lemma}
\begin{proof}
  For $k=0$ the sum becomes $q\equiv 0 \bmod q$. So suppose $k>0$. Let $\gm$ is the primitive element of $\bbF_q$. Then we have $$\sum_{\alpha\in\bbF_q}\alpha^k=\sum_{\alpha\in\bbF_q^*}\alpha^k=\sum_{j=0}^{q-2}\gm^{j\cdot k}=\sum_{j=0}^{q-2}(\gm^k)^j$$Now if $q-1\mid k$ then $\gm^k=1$ so the sum evaluates to be $-1$. If $q-1\nmid k$ then $\sum_{j=0}^{q-2}(\gm^k)^j=\frac{\gm^{k(q-1)}-1}{\gm^k-1}=0$. Hence we have the lemma. 
\end{proof}
So using this lemma we can prove if we take the sum over all evaluations of a polynomial then the sum evaluates to be zero.
\begin{lemma}{}{lem:poly-evaluate-allsum}
  Let $f\in\bbF_q[X_1,\dots,X_n]$ with $\deg(f)<n(q-1)$. Then $$\sum_{\alpha_1,\dots,\alpha_n\in\bbF_q}f(\alpha_1,\dots, \alpha_k)=0$$
\end{lemma}
\begin{proof}
  It suffices to prove this property for every monomial of $f$. So suppose $X_1^{k_1}\cdots X_n^{k_n}$ be a monomial of $f$ where $k_1+\cdots+k_n<n(q-1)$ and $k_i\in\bbZ_0$ for all $i\in[n]$. Since $k_1+\cdots+k_n<n(q-1)$ then by Pigeon Hole Principle there exists $j\in[n]$ such that $k_j< q-1$. Then
  $$\sum_{\alpha_1,\dots,\alpha_n\in\bbF_q}\alpha_1^{k_1}\cdots \alpha_n^{k_n}=\prod_{i=1}^n \lt(\sum_{\alpha_i\in\bbF_q}\alpha^{k_i}\rt)$$ Then for $j$, by \lemref{lem:power-sum} $\sum_{\alpha_j\in\bbF_q}\alpha^{k_j}=0$. So we have the lemma.
\end{proof}

The two lemmas above are the principal tools behind the Chevalley--Warning theorem. It asserts that whenever the total degree of a system is small relative to the number of variables, the number of common zeros is divisible by $p$ --- forcing, in particular, that a system with one solution must have a second.

\begin{Theorem}{Chevalley--Warning Theorem}{thm:chevalley-warning}
  \index{Chevalley-Warning theorem}%
  Let $f\in\bbF_q[X_1,\dots,X_n]$ with $\deg(f)<n$. Then the number of solutions of the equation $f(X_1,\dots, X_n)=0$ in $\bbF_q^n$ is divisible by the characteristic $p$ of $\bbF_q$. 

  Moreover, if $f(0,\dots, 0)=0$ then there exists a nontrivial solution in $\mcV(f)$ i.e. $\ov{\mcV}(f)\neq \emptyset$ or $\ov{Z}(f)\geq 1$. 
\end{Theorem}
\begin{proof}
  We first construct a polynomial and then will use \PolyEvallem. So take the polynomial $F(X_1,\dots, X_n)=1-f^{q-1}(X_1,\dots, X_n)$. By construction $F$ always takes vale $0$ or $1$. Therefore for any $\alpha_1,\dots, \alpha_n\in\bbF_q$ we have $$F(\alpha_1,\dots,\alpha_n)=1\iff f(\alpha_1,\dots, \alpha_n)=0$$So $$\sum_{\alpha_1,\dots,\alpha_n\in\bbF_q}F(\alpha_1,\dots, \alpha_k)=Z(f)\bmod p$$Now by \PolyEvallem we have $$\sum_{\alpha_1,\dots,\alpha_n\in\bbF_q}F(\alpha_1,\dots, \alpha_k)=0$$Therefore $p\mid Z(f)$. So we have the first part. 

  Now since $f(0,\dots, 0)=0$ we have $Z(f)\geq 1$. Since $p$ is a prime $p\geq 2$. So by first part $p\mid Z(f)$. Hence $Z(f)\geq 2$. So we have the second part too. 
\end{proof}

Now the bound of $\deg(f)<n$ in \CWthm is the best possible bound for the theorem to work. We can construct a polynomial in $\bbF_q[X_1,\dots, X_n]$ of degree $n$ for which the theorem doesn't hold. Let $\{\alpha_1,\dots, \alpha_n\}$ is the basis of $E=\bbF_{q^n}$ over $\bbF_q$. Consider the following polynomial $$f(X_1,\dots, X_n)=\prod_{j=0}^{n-1}\lt(\alpha_1^{q^i}X_1+\cdots +\alpha_n^{q^i}X_n\rt)=\prod_{j=0}^{n-1}\sum_{j=1}^n \alpha_j^{q^i}X_i$$Since for all $i\in\{0,1,\dots, n-1\}$, $\alpha_j^{q^i}$ are conjugates of $\alpha_j$ for all $j\in[n]$ with respect to $\bbF_q$, $f$ is a polynomial over $\bbF_q$. Suppose for any $(\gm_1,\dots, \gm_n)\in\bbF_q^n$ consider the element $\alpha_1\cdot \gm_1+\cdots+\alpha_n\cdot \gm_n\in E$. Then we evaluate $f$ on $(\gm_1,\dots, \gm_n)$ and we get $$f(\gm_1,\dots,\gm_n)=\prod_{i=0}^{n-1}\sum_{j=1}^n \alpha_j^{q^i}\cdot \gm_j=\prod_{i=0}^{n-1}\lt(\sum_{j=1}^n \alpha_j\cdot \gm_j\rt)^{q^i}=N_{E/\bbF_q}(\alpha_1\cdot \gm_1+\cdots+\alpha_n\cdot \gm_n)$$Hence $f(\gm_1,\dots,\gm)=0$ if and only if $N_{E/\bbF_q}(\alpha_1\cdot \gm_1+\cdots+\alpha_n\cdot \gm_n)=0$ which holds only for $\alpha_1\cdot \gm_1+\cdots+\alpha_n\cdot \gm_n=0$. Since  $\{\alpha_1,\dots, \alpha_n\}$ forms a basis of $E$ over $\bbF_q$ we have $f(\gm_1,\dots,\gm)=0$ if and only if $\gm_j=0$ for all $j\in[n]$. Hence $Z(f)=1$. The \CWthm clearly doesn't hold in this case. 

Now we can extend the \CWthm to a system of polynomial equations $f_1=0,\dots,f_k=0$ with $f_1,\dots,f_k\in\bbF_q[X_1,\dots,X_n]$. In this case we are interested in the number of common solutions. 
\begin{Theorem}{}{}
  Let $f_1,\dots, f_k\in\bbF_q[X_1,\dots, X_n]$ with $\deg(f_1\cdots f_k)=\deg(f_1)+\cdots+\deg(f_k)<n$. Then $p\mid Z(f_1,\dots, f_k)$ where $p=\Char(\bbF_q)$. 

  Moreover, if $(0,\dots, 0)\in \mcV(f_1,\dots, f_k)$ then $\ov{\mcV}(f_1,\dots, f_k)\neq\emptyset$ i.e. $\ovZ(f_1,\dots, f_k)\geq 1$. 
\end{Theorem}
\begin{proof}
  Again we construct a polynomial. Consider the polynomial $F\in\bbF_q[X_1,\dots, X_n]$ such that $F=(1-f_1^{q_1})\cdots (1-f_k^{q-1})$. Therefore for any $\alpha_1,\dots,\alpha_n\in\bbF_q$ we have $F(\alpha_1,\dots,\alpha_n)=1$ if and only if $f_i(\alpha_1,\dots, \alpha_n)=0$ for all $i\in[k]$. Now $\deg(F)\leq \deg(f_1)+\cdots+\deg(f_k)<n(q-1)$. Now by applying \CWthm for $F$ we get the theorem. 
\end{proof}

\subsection{Lower Bounds on the Number of Solutions}

In the next theorem we will show if we have a system of polynomial equations $f_1=0,\dots,f_k=0$ with $f_1,\dots,f_k\in\bbF_q[X_1,\dots,X_n]$ then for any affine subspace $W$ all its parallel shifts have same number of solutions modulo the characteristic of $\bbF_q$. If $W_1$ and $W_2$ are two affine subspaces of $\bbF_q^n$ of same dimension then we call they are parallel affine subspaces if they are obtained by translation from the same linear subspace. 
\begin{lemma}{}{lem:solutions-affine-subspace-parallel-shift}
  Let $f_1,\dots, f_k\in\bbF_q[X_1,\dots, X_n]$.  If $W_1$ and $W_2$ are two parallel affine subspaces of $\bbF_q^n$ of dimension $d$ where $d=\deg(f_1)+\cdots+\deg(f_k)$ then $$|W_1\cap \mcV(f_1,\dots,f_k)|\equiv |W_2\cap \mcV(f_1,\dots, f_k)|\bmod p$$where $p=\Char(\bbF_q)$. 
\end{lemma}
\begin{proof}
  Now if $W_1=W_2$, we get $|W_1\cap \mcV(f_1,\dots,f_k)|=|W_2\cap \mcV(f_1,\dots, f_k)|$. So we get the theorem. So we can assume $W_1\neq W_2$. Now by change of coordinates we can take \begin{align*}
    & W_1 = \{(\alpha_1,\dots, \alpha_n)\in\bbF_q^n\colon \alpha_i=0,\ \forall\ i\in[n]\}\\ 
    & W_2 = \{(\alpha_1,\dots, \alpha_n)\in\bbF_q^n\colon \alpha_1=1\text{ and }\alpha_i=0,\ \forall\ 2\leq i\leq n\}\\ 
  \end{align*}
  Again we construct a polynomial $G\in\bbF_q[X_1,\dots, X_n]$ where $$G(X_1,\dots, X_n)=(-1)^{n-d}(X_1^{q-2}+\cdots X_1+1)(X_2^{q-1}-1)\cdots (X_n^{q-1}-1)$$So $\deg(G)=(n-d)(q-1)-1$. 
  \begin{observation}
    With the above construction of $G$ for any $(\alpha_1,\dots,\alpha_n)\in\bbF_q^n$ \begin{enumerate}[label=(\roman*)]
      \item if $(\alpha_1,\dots,\alpha_n)\in W_1$ then $G(\alpha_1,\dots, \alpha_n)=-1$
      \item if $(\alpha_1,\dots,\alpha_n)\in W_2$ then $G(\alpha_1,\dots, \alpha_n)=1$
      \item if $(\alpha_1,\dots,\alpha_n)\notin W_1\cup W_2$ then $G(\alpha_1,\dots, \alpha_n)=0$
    \end{enumerate}
  \end{observation}
  Then we construct the following polynomial $F\in\bbF_q[X_1,\dots, X_n]$ such that $$F=(1-f_1^{q-1})\cdots (1-f_k^{q-1})G$$So now we have $\deg(F)\leq d(q-1)+(n-d)(q-1)-1=n(q-1)-1$. Therefore for any $(\alpha_1,\dots,\alpha_n)\in\bbF_q^n$, if $(\alpha_1,\dots,\alpha_n)\in W_1\cap \mcV(f_1,\dots, f_k)$ then $F(\alpha_1,\dots, \alpha_n)=1$. If $(\alpha_1,\dots,\alpha_n)\in W_2\cap \mcV(f_1,\dots, f_k)$ then $F(\alpha_1,\dots, \alpha_n)=-1$ and $0$ elsewhere. So $$\sum_{\alpha_1,\dots,\alpha_n\in\bbF_q}F(\alpha_1,\dots,\alpha_n)=|W_1\cap \mcV(f_1,\dots,f_k)|- |W_2\cap \mcV(f_1,\dots, f_k)|$$Now by using \CWthm on $F$ we have the lemma. 
\end{proof}

The \AffSubspaceShiftlem above can be used to obtain a lower bound on $Z(f_1,\dots,f_k)$: the \CWthm guarantees that $p\mid Z(f_1,\dots,f_k)$ whenever $d<n$, but by counting how many parallel affine hyperplanes each must contain a solution we can say much more.
\begin{Theorem}{}{}
  Let $f_1,\dots, f_k\in\bbF_q[X_1,\dots, X_n]$ with $d=\deg(f_1)+\cdots+\deg(f_k)<n$. If $Z(f_1,\dots, f_k)\geq 1$ then $Z(f_1,\dots, f_k)\geq q^{n-d}$
\end{Theorem}
\begin{proof}
  We will prove this by analyzing two cases. 
  \paragraph*{Case I:} Let there exists an affine subspace $W_1$ of $\bbF_q^n$ with dimension $d$ such that $|W_1\cap \mcV(f_1,\dots, f_k)|\not\equiv 0\bmod p$. Then by \AffSubspaceShiftlem for all parallel affine spaces $W_2$ of $W_1$ we have $|W_2\cap \mcV(f_1,\dots,f_k)|\not\equiv 0\bmod p$. In particular $|W_2\cap \mcV(f_1,\dots, f_k)|\geq 1$. Since $W_1$ is of dimension $d$ there are $q^{n-d}$ distinct affine subspaces parallel to $W_1$. Therefore 
  $$\mcV(f_1,\dots, f_k)=\bigsqcup_{\substack{W\colon\text{affine subspace}\\ \text{parallel to $W_1$}}}\Big(W\cap \mcV(f_1,\dots,f_k)\Big)$$
  For each such $W$ we have $|W\cap\mcV(f_1,\dots,f_k)|\geq 1$. Therefore 
  $$Z(f_1,\dots, f_k)=\sum_{\substack{W\colon\text{affine subspace}\\ \text{parallel to $W_1$}}}|W\cap\mcV(f_1,\dots,f_k)|\geq q^{n-d}$$
  Therefore we have the theorem in this case. 
  \paragraph*{Case II:} We have $|W\cap\mcV(f_1,\dots,f_k)|\equiv 0\bmod p$. for all affine subspaces $W$ of $\bbF_q^n$ of dimension $d$. Since $Z(f_1,\dots,f_k)\geq 1$ there exists $ k\in[d]$ such that for any affine subspace $W'$ of dimension $k$, $|W'\cap \mcV(f_1,\dots,f_k)|\equiv 0\bmod p$ but there is an affine subspace $W''$ of dimension $(k-1)$ such that $|W''\cap \mcV(f_1,\dots, f_k)|\not\equiv 0\bmod p$. Now consider all affine subspaces $W'$ of dimension $k$ containing $W''$. There are exactly $\frac{q^{n-(k-1)}-1}{q-1}$ many such affine subspaces. For each such affine subspace $W'$ we have $$|(W'\setminus W'')\cap \mcV(f_1,\dots,f_k)|=|W'\cap \mcV(f_1,\dots,f_k)|-|W''\cap \mcV(f_1,\dots,f_k)|\not\equiv0\bmod p$$So $|(W'\setminus W'')\cap \mcV(f_1,\dots,f_k)|\geq 1$. Therefore we have $$\mcV(f_1,\dots, f_k)=\Big(W''\cap \mcV(f_1,\dots,f_k)\Big)\sqcup\bigsqcup_{\substack{W'\colon \text{affine subspace}\\ \text{of dimension $k$}\\ \text{contains $W''$}}}\Big((W'\setminus W'')\cap \mcV(f_1,\dots,f_k)\Big)$$Hence we have $$Z(f_1,\dots, f_k)=|W''\cap \mcV(f_1,\dots,f_k)|+\sum_{\substack{W'\colon \text{affine subspace}\\ \text{of dimension $k$}\\ \text{contains $W''$}}}|(W'\setminus W'')\cap \mcV(f_1,\dots,f_k)|\geq 1+\frac{q^{n-(k-1)}-1}{q-1}>q^{n-d}$$So we have the theorem. 
\end{proof}

The lower bound on the number of solutions above is best possible even for $k=1$ i.e. for any positive integers $d$ and $n$ with $d<n$ there is a polynomial $f\in\bbF_q[X_1,\dots, X_n]$ of degree $d$ such that $f(X_1,\dots, X_n)=0$ has exactly $q^{n-d}$ solutions in $\bbF_q^n$. Consider the polynomial defined in the support \CWthm, $\prod_{i=0}^{n-1}\sum_{j=1}^n\alpha_j^{q^i}X_j$ where $\{\alpha_j\colon j\in[n]\}$ was a basis of $\bbF_{q^n}$ over $\bbF_q$. But in this example we take the same polynomial but on $d$ variables but the solution space is over $\bbF_{q^n}$. Hence our polynomial in this case is $f(X_1,\dots,X_n)=g(X_1,\dots,X_d)=\prod_{i=0}^{n-1}\sum_{j=1}^d\alpha_j^{q^i}X_j$. Since $f$ had only one solution, $g$ has a solution at $(\gm_1,\dots,\gm_n)$ if $\gm_i=0$ for all $i\in[d]$. Therefore $g$ has exactly $q^{n-d}$ solutions in $\bbF_q^n$. 

\subsection{Upper Bounds on the Number of Solutions}

We now give upper bounds on the number of solutions of a polynomial equation. The first is the celebrated Schwartz--Zippel lemma, which bounds the probability of a random point being a root in terms of total degree.
\begin{lemma}{Schwartz-Zippel Lemma}{lem:schwartz-zippel}
  Let $f\in\bbF_q[X_1,\dots, X_n]$ with $\deg(f)=d$ be a non-zero polynomial. If $S\subseteq \bbF_q$ be any subset then $$\Pr[\alpha_1,\dots,\alpha_n\in S] [f(\alpha_1,\dots,\alpha_n)=0]\leq \frac{d}{|S|}$$
\end{lemma}
This lemma immediately gives an upper bound of $d\cdot q^{n-1}$ solutions for the equations $f(X_1,\dots, X_n)=0$ in $\bbF_q^n$. Hence 
\begin{corolary}{}{}
  Let $f\in\bbF_q[X_1,\dots, X_n]$ with $\deg(f)=d$ be a non-zero polynomial. Then the equation $f(X_1,\dots, X_n)=0$ has at most $d\cdot q^{n-1}$ solutions in $\bbF_q^n$. 
\end{corolary}

We can give another probabilistic upper bound when we control the degree of $f$ in each variable separately rather than just the total degree.
\begin{Theorem}{Individual Degree Bound}{thm:individual-degree-bound}
  Let $f\in\bbF_q[X_1,\dots,X_n]$ be a non-zero polynomial with $\deg_{X_i}(f)\leq d$ for all $i\in[n]$. Then for $(\alpha_1,\dots,\alpha_n)$ chosen uniformly at random from $\bbF_q^n$,
  $$\Pr[\alpha_1,\dots,\alpha_n\in\bbF_q][f(\alpha_1,\dots,\alpha_n)=0]\leq 1-\lt(1-\frac{d}{q}\rt)^{n}$$
\end{Theorem}
\begin{proof}
  We proceed by induction on $n$. For the base case $n=1$, since $f\in\bbF_q[X_1]$ is non-zero with $\deg_{X_1}(f)\leq d$, it has at most $d$ roots in $\bbF_q$. Hence $\Pr[x_1\in\bbF_q][f(x_1)=0]\leq d/q=1-(1-d/q)^1$. 
  
  Assume the bound holds for $n-1$ variables. Let $t=\deg_{X_n}(f)$. Then write $f$ as a polynomial in $X_n$ with coefficients in $\bbF_q[X_1,\dots,X_{n-1}]$:
  $$f(X_1,\dots,X_n)=\sum_{i=0}^{t}f_i(X_1,\dots,X_{n-1})\cdot X_n^k.$$
  Since $\deg_{X_j}(f)\leq d$ for each $j\in[n-1]$, every $f_i$ satisfies $\deg_{X_j}(f_i)\leq d$.

  Let $\alpha_1,\dots,\alpha_{n-1}$ be uniformly random points on $\bbF_q$ and $\alpha_n$ be uniform on $\bbF_q$, independently. Then by induction hypothesis $$\Pr[\alpha_1,\dots,\alpha_{n-1}\in\bbF_q] [f_t(\alpha_1,\dots,\alpha_{n-1})=0]\leq 1-\lt(1-\frac{d}{q}\rt)^{n-1}$$Now we have \begin{align*}
    \underset{\alpha_1,\dots,\alpha_n\in\bbF_q}{\bbP} [f(\alpha_1,\dots,\alpha_n)=0] &\leq \underset{\alpha_1,\dots,\alpha_{n-1}\in\bbF_q}{\bbP} [f_t(\alpha_1,\dots,\alpha_{n-1})=0]+ \underset{\alpha_1,\dots,\alpha_n\in\bbF_q}{\bbP} [f(\alpha_1,\dots, \alpha_{n-1},\alpha_n)=0\wedge f_t(\alpha_1,\dots,\alpha_{n-1})\neq 0]\\ 
    & \leq \underset{\alpha_1,\dots,\alpha_{n-1}\in\bbF_q}{\bbP} [f_t(\alpha_1,\dots,\alpha_{n-1})=0]\;+ \\ 
    &\qquad\qquad \underset{\alpha_n\in\bbF_q}{\bbP} [f(\alpha_1,\dots, \alpha_{n-1},\alpha_n)=0\mid f_t(\alpha_1,\dots,\alpha_{n-1})\neq 0]\underset{\alpha_1,\dots,\alpha_{n-1}\in\bbF_q}{\bbP}[f_t(\alpha_1,\dots,\alpha_{n-1})\neq 0]
  \end{align*}
  Let $p\coloneqq \underset{\alpha_1,\dots,\alpha_{n-1}\in\bbF_q}{\bbP} [f_t(\alpha_1,\dots,\alpha_{n-1})=0]$. For the second term after fixing $\alpha_1,\dots,\alpha_{n-1}$ we have $\underset{\alpha_n\in\bbF_q}{\bbP} [f(\alpha_1,\dots, \alpha_{n-1},\alpha_n)=0\mid f_t(\alpha_1,\dots,\alpha_{n-1})\neq 0]\leq\frac{d}q$. Now $$\underset{\alpha_1,\dots,\alpha_{n-1}\in\bbF_q}{\bbP}[f_t(\alpha_1,\dots,\alpha_{n-1})\neq 0]=1-\underset{\alpha_1,\dots,\alpha_{n-1}\in\bbF_q}{\bbP}[f_t(\alpha_1,\dots,\alpha_{n-1})= 0]$$ Therefore together we get $$\underset{\alpha_1,\dots,\alpha_n\in\bbF_q}{\bbP} [f(\alpha_1,\dots,\alpha_n)=0]\leq p+(1-p)\frac{d}{q}=\frac{d}{q}+\lt(1-\frac{d}{q}\rt)p$$By induction hypothesis $p\leq 1-\lt(1-\frac{d}q\rt)^{n-1}$. Therefore we have $$\underset{\alpha_1,\dots,\alpha_n\in\bbF_q}{\bbP} [f(\alpha_1,\dots,\alpha_n)=0]\leq \frac{d}{q}+\lt(1-\frac{d}{q}\rt)\lt[1-\lt(1-\frac{d}{q}\rt)^{n-1}\rt]=1-\lt(1-\frac{d}{q}\rt)^{n}$$So we have the lemma. 
\end{proof}

We can give another upper bound based on individual degrees of each variables. If degree of each variable is at most $d$ then there can be at most $d(q^{n-1}-1)$ many non-trivial solutions. For this we consider $f$ to be \emph{homogeneous} i.e. total degree of each monomial is same. 
\begin{Theorem}{}{}
  Let $f\in\bbF_q[X_1,\dots, X_n]$ be a homogeneous with $\deg(f)=d\geq 1$. Then the equation $f(X_1,\dots,X_n)=0$ has at most $d(q^{n-1}-1)$ non-trivial solutions. 
\end{Theorem}
\begin{proof}
  We use double induction on $n$ and $d$. For the base case $n=1$ then $f=aX_1^d$ for some $a\neq 0$, so the only root is $X_1=0$, giving $\ov{Z}(f)=0\leq d(q^0-1)=0$. If $d=1$ then $f=a_1X_1+\cdots+a_nX_n$ is a non-zero linear form whose zero set is an $(n-1)$-dimensional hyperplane, so $\ov{Z}(f)=q^{n-1}-1=d(q^{n-1}-1)$.

  Suppose $n>1$, $d>1$, and the bound holds for non-constant homogeneous polynomials in at most $n$ variables of degree less than $d$ and for those in fewer than $n$ variables of degree at most $d$. We consider two cases.

  \paragraph*{Case 1: $X_1\mid f$.} Write $f=X_1\cdot g$ where $g\in\bbF_q[X_1,\dots,X_n]$ is homogeneous of degree $d-1$. Every nontrivial solution $(\alpha_1,\dots,\alpha_n)\in\ovZ(f)$ with $\alpha_i\in\bbF_q$ for all $i\in[n]$, satisfies $\alpha_1=0$ or $g(\alpha_1,\dots,\alpha_n)=0$. Solutions with $\alpha_1=0$: any $(0,\alpha_2,\dots,\alpha_n)\neq 0$ works, giving $q^{n-1}-1$ solutions. Solutions with $\alpha_1\neq 0$: these are nontrivial zeros of $g$, and by the induction hypothesis on $d$ there are at most $(d-1)(q^{n-1}-1)$. Hence $$\ov{Z}(f)\leq (q^{n-1}-1)+(d-1)(q^{n-1}-1)=d(q^{n-1}-1).$$

  \paragraph*{Case 2: $X_1\nmid f$.} Write $f=\sum_{k=0}^{d}f_k(X_2,\dots,X_n)\cdot X_1^k$ with $f_k$ homogeneous of degree $d-k$. Since $X_1\nmid f$ we have $f_0\neq 0$, so for every $\alpha\in\bbF_q$ such that $f(\alpha,X_2,\dots,X_n)$ is a non-zero polynomial of degree $d$ in $n-1$ variables we have by \SZlem it has at most $d\cdot q^{n-2}$ zeros in $\bbF_q^{n-1}$. Now summing over $\alpha\in\bbF_q^*$ gives at most $(q-1)\cdot d\cdot q^{n-2}$ nontrivial solutions with $\alpha_1\neq 0$. For solutions with $\alpha_1=0$: $f(0,X_2,\dots,X_n)=f_0(X_2,\dots,X_n)$ is a nonzero homogeneous polynomial in $n-1$ variables of degree $d$, so by the induction hypothesis on $n$ it has at most $d(q^{n-2}-1)$ nontrivial zeros. Altogether, $\ov{Z}(f)\leq (q-1)d\cdot q^{n-2}+d(q^{n-2}-1)=d\lt(q^{n-1}-q^{n-2}+q^{n-2}-1\rt)=d(q^{n-1}-1)$.
\end{proof}
\subsection{Average Number of Solutions}
We now investigate the \emph{average number of solutions} of a polynomial equations. Let $d\in\bbN$. Then define $\Om_d\coloneqq \bbF_q^{\leq d}[X_1,\dots, X_n]$ and $\om(d)$ be the set of $n-$tuples $(i_1,\dots,i_n)\in\bbZ_0$ such that $i_1+\cdots +i_n\leq d$. Then 
\begin{observation}\label{obs:Omd-omd}
  With the definitions of $\Om_d$ and $\om(d)$ as above we have $|\Om_d|=q^{|\om(d)|}$. 
\end{observation}
\begin{lemma}{}{lem:average-solutions}
  With the notations of $\Om_d$ and $\om(d)$ as defined above we have $$\frac1{|\Om_d|}\sum_{f\in\Om_d}Z(f)=q^{n-1}$$
\end{lemma}
\begin{proof}
  We have $$\sum_{f\in\Om_d}Z(f)=\sum_{f\in\Om_d}\sum_{\substack{\alpha_1,\dots,\alpha_n\in\bbF_q\\ f(\alpha_1,\dots,\alpha_n)=0}}1=\sum_{\alpha_1,\dots,\alpha_n\in\bbF_q}\sum_{\substack{f\in\Om_d\\ f(\alpha_1,\dots,\alpha_n)=0}}1$$Let for any $f\in\Om_d$ $$f(X_1,\dots,X_n)=\sum_{(i_1,\dots,i_n)\in\om(d)}f_{i_1,\dots,i_n}X_1^{i_1}\cdots X_n^{i_n}\quad f_{i_1,\dots,i_n}\in\bbF_q,\ \forall\ (i_1,\dots, i_n)\in\om(d)$$Then for a fixed $\alpha_1,\dots, \alpha_n\in\bbF_q$ we get to choose the coefficients $f_{i_1,\dots, i_n}$ for all $(i_1,\dots,i_n)\in\om(d)\setminus\{(0,\dots, 0)\}$ arbitrarily and then determining the constant term $f_{0,\dots, o}$ to make $f(\alpha_1,\dots,\alpha_n)=0$. Thus the number of $f\in\om_d$ with $f(\alpha_1,\dots, \alpha_n)=0$ is equal to $q^{|\om(d)|-1}$. Therefore $\sum_{f\in\Om_d}Z(f)=q^n\cdot q^{|\om(d)|-1}=|\Om_d|\cdot q^{n-1}$. So we have the lemma. 
\end{proof}
The proof of the above lemma gives us the following observation:
\begin{observation}\label{obs:num-poly-zero-at-point}
  For any $\alpha_1,\dots,\alpha_n\in\bbF_q$, the number of polynomials $f\in\Om_d$ such that $f(\alpha_1,\dots,\alpha_n)=0$ is $q^{|\om(d)|-1}$. 
\end{observation}

The above lemma implies that a polynomial equation in $n$ variables has on average $q^{n-1}$ solutions in $\bbF_q^n$. So we now consider average deviation from the average number of solutions, $q^{n-1}$. 
\begin{Theorem}{}{}
  With the notations of $\Om_d$ and $\om(d)$ as defined above we have
  $$\frac1{|\Om_d|}\sum_{f\in\Om_d}\Big(Z(f)-q^{n-1}\Big)^2=q^{n-1}-q^{n-2}$$
\end{Theorem}
\begin{proof}
  Let $\alpha=(\alpha_1,\dots,\alpha_n),\beta=(\beta_1,\dots,\beta_n)\in\bbF_q^n$ be any two points. Then we get
  $$\sum_{f\in\Om_d}Z(f)^2  = \sum_{f\in\Om_d}\lt( \sum_{\substack{\alpha\in\bbF_q^n\\ f(\alpha)=0}} 1 \rt)^2= \sum_{f\in\Om_d}\sum_{\substack{\alpha\in\bbF_q^n\\ f(\alpha)=0}}\sum_{\substack{\beta\in\bbF_q^n\\ f(\beta)=0}}1=\sum_{\alpha,\beta\in\bbF_q^n}\sum_{\substack{f\in\Om_d\\  f(\alpha)=f(\beta)=0}}1$$
  Now if $\alpha=\beta$ then by \obsref{obs:num-poly-zero-at-point} we have $\sum_{f\in\Om_d,\   f(\alpha)=0}1=q^{|\om(d)|-1}$. Now suppose $\alpha\neq\beta$. This gives two linear equations for the coefficients of $f$. Therefore there are $q^{|\om(d)|-2}$ polynomials $f\in\Om_d$ such that $f(\alpha)=f(\beta)=0$. Hence we obtain \begin{align*}
    \sum_{f\in\Om_d}Z(f)^2  =  & = \sum_{\alpha\in\bbF_q^n}|\Om_d|\cdot q^{-1}+\sum_{\alpha,\beta\in\bbF_q^n,\alpha\neq \beta}|\Om_d|\cdot q^{-2} \byo{obs:Omd-omd}\\ 
    & = q^n\cdot |\Om_d|\cdot q^{-1}+q^n(q^n-1)|\Om_d|\cdot q^{-2}\\ 
    & = |\Om_d|\Big( q^{n-1}+q^{2n-2}-q^{n-2} \Big)
  \end{align*}
  So now we get \begin{align*}
    \sum_{f\in\Om_d}(Z(f)-q^{n-1})^2 & = \sum_{f\in\Om_d}Z(f)^2 + 2\cdot q^{n-1}\sum_{f\in\Om_d}Z(f)+q^{2n-2}\sum_{f\in\Om_d}1\\ 
    & = |\Om_d|\Big(q^{2n-2}+q^{n-1}-q^{n-2}\Big)-2\cdot q^{n-1}\cdot |\Om_d|\cdot q^{n-1}+q^{2n-2}\cdot |\Om_d|\\ 
    & = |\Om_d|(q^{n-1}-q^{n-2})
  \end{align*}
  So we have the result. 
\end{proof}
From this result we can expect that $|Z(f)-q^{n-1}|$ is often $O\lt(q^{\nicefrac{(q-1)}{2}}\rt)$. Later we will see instances of such expected behavior for various polynomial equations. In the next sections we will talk about character sums over polynomial evaluations. 
